\subsection{Structural Abstraction of the Unicode Standard}

Python's \texttt{str} object represents text. More precisely, it represents a sequence of Unicode codepoints. This is already a higher-level structure than a C-style byte array.

A Python string does not expose ``the bytes of the text'' during normal string operations. It exposes text components. When bytes are needed, the string must be encoded.

\begin{verbatim}
text = "D'oh!"
raw = text.encode("utf-8")

print(f"text: {text!r}")
print(f"raw:  {raw!r}")

print()
print(f"type(text): {type(text)}")
print(f"type(raw):  {type(raw)}")
\end{verbatim}

Possible output:

\begin{verbatim}
text: "D'oh!"
raw:  b"D'oh!"

type(text): <class 'str'>
type(raw):  <class 'bytes'>
\end{verbatim}

The visible result is simple, but the layers are different:

\begin{itemize}
    \item \texttt{str}: text-level object;
    \item Unicode codepoint: abstract character number;
    \item glyph: visible drawing on the screen;
    \item \texttt{bytes}: transmitted or stored byte values.
\end{itemize}

The following inspection shows that \texttt{str} and \texttt{bytes} are different object types with different available operations.

\begin{verbatim}
text = "Ni!"
raw = b"Ni!"

selected_names = [
    "encode",
    "decode",
    "__len__",
    "__getitem__",
    "split",
    "join",
    "lower",
    "hex",
]

for obj_name, obj in [("text", text), ("raw", raw)]:
    print(f"{obj_name}: {type(obj)}")
    for name in selected_names:
        print(f"  {name:<12} {name in dir(obj)}")
    print()
\end{verbatim}

Possible output:

\begin{verbatim}
text: <class 'str'>
  encode       True
  decode       False
  __len__      True
  __getitem__  True
  split        True
  join         True
  lower        True
  hex          False

raw: <class 'bytes'>
  encode       False
  decode       True
  __len__      True
  __getitem__  True
  split        True
  join         True
  lower        True
  hex          True
\end{verbatim}

The important boundary is directional:

\begin{itemize}
    \item \texttt{str.encode(...)} produces \texttt{bytes};
    \item \texttt{bytes.decode(...)} produces \texttt{str}.
\end{itemize}

\subsubsection{Distinction Between Abstract Codepoints, Visible Glyphs, and Transmitted Byte Serializations}

A Unicode codepoint is an abstract number assigned to a text character. Python exposes that number with \texttt{ord()}.

\begin{verbatim}
chars = ["A", "ă", "π", "☃", "😀"]

for ch in chars:
    print(f"ch={ch!r:<4} ord={ord(ch):<7} unicode=U+{ord(ch):06X}")
\end{verbatim}

Possible output:

\begin{verbatim}
ch='A'  ord=65      unicode=U+000041
ch='ă'  ord=259     unicode=U+000103
ch='π'  ord=960     unicode=U+0003C0
ch='☃'  ord=9731    unicode=U+002603
ch='😀' ord=128512  unicode=U+01F600
\end{verbatim}

This is not byte inspection. It is text inspection. The codepoint tells us which abstract Unicode character is present.

The byte serialization is obtained only after encoding.

\begin{verbatim}
chars = ["A", "ă", "π", "☃", "😀"]

for ch in chars:
    raw = ch.encode("utf-8")
    print(f"ch={ch!r:<4} codepoint=U+{ord(ch):06X} utf8={raw!r}")
\end{verbatim}

Possible output:

\begin{verbatim}
ch='A'  codepoint=U+000041 utf8=b'A'
ch='ă'  codepoint=U+000103 utf8=b'\xc4\x83'
ch='π'  codepoint=U+0003C0 utf8=b'\xcf\x80'
ch='☃'  codepoint=U+002603 utf8=b'\xe2\x98\x83'
ch='😀' codepoint=U+01F600 utf8=b'\xf0\x9f\x98\x80'
\end{verbatim}

The same one-character string can need one, two, three, or four bytes in UTF-8.

\begin{verbatim}
chars = ["A", "ă", "π", "☃", "😀"]

print(f"{'char':<6} {'len(str)':<10} {'len(utf8)':<10} {'utf8 bytes'}")

for ch in chars:
    raw = ch.encode("utf-8")
    print(f"{ch!r:<6} {len(ch):<10} {len(raw):<10} {raw!r}")
\end{verbatim}

Possible output:

\begin{verbatim}
char   len(str)   len(utf8)  utf8 bytes
'A'    1          1          b'A'
'ă'    1          2          b'\xc4\x83'
'π'    1          2          b'\xcf\x80'
'☃'    1          3          b'\xe2\x98\x83'
'😀'   1          4          b'\xf0\x9f\x98\x80'
\end{verbatim}

So \texttt{len(ch)} counts Python string components, not bytes.

There is another layer: the visible glyph. A glyph is what appears on the screen. A visible glyph can sometimes be built from more than one Unicode codepoint.

\begin{verbatim}
text1 = "é"
text2 = "e\u0301"

print(f"text1: {text1!r}")
print(f"text2: {text2!r}")

print()
print(f"printed text1: {text1}")
print(f"printed text2: {text2}")

print()
print(f"len(text1): {len(text1)}")
print(f"len(text2): {len(text2)}")
\end{verbatim}

Possible output:

\begin{verbatim}
text1: 'é'
text2: 'é'

printed text1: é
printed text2: é

len(text1): 1
len(text2): 2
\end{verbatim}

The two strings may look nearly identical when printed. Internally, they are not the same sequence of codepoints.

\begin{verbatim}
text1 = "é"
text2 = "e\u0301"

print("text1 components")
for ch in text1:
    print(f"{ch!r:<4} U+{ord(ch):04X}")

print()
print("text2 components")
for ch in text2:
    print(f"{ch!r:<4} U+{ord(ch):04X}")

print()
print(f"text1 == text2: {text1 == text2}")
\end{verbatim}

Possible output:

\begin{verbatim}
text1 components
'é'  U+00E9

text2 components
'e'  U+0065
'́'  U+0301

text1 == text2: False
\end{verbatim}

The second string contains:

\begin{itemize}
    \item \texttt{U+0065}: the letter \texttt{e};
    \item \texttt{U+0301}: a combining acute accent.
\end{itemize}

A combining mark modifies the visual rendering of the preceding character, but it is still a separate codepoint inside the Python string.

The standard library module \texttt{unicodedata} can show the names of those components.

\begin{verbatim}
import unicodedata

text = "e\u0301"

for ch in text:
    name = unicodedata.name(ch)
    print(f"{ch!r:<4} U+{ord(ch):04X} {name}")
\end{verbatim}

Possible output:

\begin{verbatim}
'e'  U+0065 LATIN SMALL LETTER E
'́'  U+0301 COMBINING ACUTE ACCENT
\end{verbatim}

Normalization can convert between some equivalent Unicode spellings.

\begin{verbatim}
import unicodedata

text1 = "é"
text2 = "e\u0301"

n1 = unicodedata.normalize("NFC", text1)
n2 = unicodedata.normalize("NFC", text2)

print(f"text1 == text2: {text1 == text2}")
print(f"n1 == n2:       {n1 == n2}")

print()
print(f"n1: {n1!r} len={len(n1)}")
print(f"n2: {n2!r} len={len(n2)}")
\end{verbatim}

Possible output:

\begin{verbatim}
text1 == text2: False
n1 == n2:       True

n1: 'é' len=1
n2: 'é' len=1
\end{verbatim}

This example separates three notions:

\begin{itemize}
    \item codepoint sequence: what Python's \texttt{str} stores and indexes;
    \item glyph rendering: what the terminal, editor, browser, or font draws;
    \item byte serialization: what encoding produces for storage or transmission.
\end{itemize}

A Python string is therefore not a drawing and not a byte buffer. It is a text sequence.

\subsubsection{Text vs. Bytes: Why \texttt{str} Stores Text and \texttt{bytes} Stores Raw Byte Values}

The \texttt{str} type stores text-level components. The \texttt{bytes} type stores raw byte values. A byte value is an integer from \texttt{0} to \texttt{255}.

The difference appears immediately during indexing.

\begin{verbatim}
text = "Aă"
raw = text.encode("utf-8")

print(f"text: {text!r}")
print(f"raw:  {raw!r}")

print()
print(f"text[0]: {text[0]!r} type={type(text[0])}")
print(f"raw[0]:  {raw[0]!r} type={type(raw[0])}")
\end{verbatim}

Possible output:

\begin{verbatim}
text: 'Aă'
raw:  b'A\xc4\x83'

text[0]: 'A' type=<class 'str'>
raw[0]:  65 type=<class 'int'>
\end{verbatim}

Indexing a string returns a one-character string. Indexing a bytes object returns an integer byte value.

Iteration shows the same split.

\begin{verbatim}
text = "Aă"
raw = text.encode("utf-8")

print("string iteration")
for ch in text:
    print(f"{ch!r:<4} U+{ord(ch):04X}")

print()
print("bytes iteration")
for byte in raw:
    print(f"{byte:<3} 0x{byte:02X}")
\end{verbatim}

Possible output:

\begin{verbatim}
string iteration
'A'  U+0041
'ă'  U+0103

bytes iteration
65  0x41
196 0xC4
131 0x83
\end{verbatim}

The string has two text components. The UTF-8 byte serialization has three byte values.

\begin{verbatim}
text = "Aă"
raw = text.encode("utf-8")

print(f"len(text): {len(text)}")
print(f"len(raw):  {len(raw)}")
\end{verbatim}

Possible output:

\begin{verbatim}
len(text): 2
len(raw):  3
\end{verbatim}

This difference explains why Python does not automatically mix \texttt{str} and \texttt{bytes} in ordinary operations.

\begin{verbatim}
text = "Ni!"
raw = b"Ni!"

print(f"text == raw: {text == raw}")

try:
    result = text + raw
    print(result)
except TypeError as err:
    print("concatenation failed")
    print(f"message: {err}")
\end{verbatim}

Possible output:

\begin{verbatim}
text == raw: False
concatenation failed
message: can only concatenate str (not "bytes") to str
\end{verbatim}

The values look similar to a human reader, but they are not the same layer. Python refuses to guess which encoding should be used.

The correct conversion must be explicit.

\begin{verbatim}
text = "No soup for you!"

raw = text.encode("utf-8")
again = raw.decode("utf-8")

print(f"text:  {text!r}")
print(f"raw:   {raw!r}")
print(f"again: {again!r}")

print()
print(f"text == again: {text == again}")
\end{verbatim}

Possible output:

\begin{verbatim}
text:  'No soup for you!'
raw:   b'No soup for you!'
again: 'No soup for you!'

text == again: True
\end{verbatim}

When data crosses a file, socket, database driver, HTTP body, or binary protocol boundary, bytes usually appear. When the program wants to reason about human text, it should decode those bytes into \texttt{str}.

\begin{verbatim}
incoming = b"Chuck Norris counted to 42."

text = incoming.decode("utf-8")

print(f"incoming: {incoming!r}")
print(f"text:     {text!r}")

print()
print(f"type(incoming): {type(incoming)}")
print(f"type(text):     {type(text)}")
\end{verbatim}

Possible output:

\begin{verbatim}
incoming: b'Chuck Norris counted to 42.'
text:     'Chuck Norris counted to 42.'

type(incoming): <class 'bytes'>
type(text):     <class 'str'>
\end{verbatim}

The practical rule is:

\begin{quote}
Use \texttt{str} for text meaning. Use \texttt{bytes} for raw byte values. Move between them only with an explicit encoding or decoding operation.
\end{quote}